El presente trabajo tomó como punto de partida un proyecto desarrollado previamente sobre un brazo robótico cartesiano destinado a aplicaciones de manipulación. Durante la revisión del sistema existente, se identificaron deficiencias mecánicas que afectaban el desplazamiento de los tres ejes. Estas fallas limitaban la fluidez, precisión y confiabilidad del movimiento, lo que reducía la capacidad del equipo para operar de forma estable dentro de una aplicación productiva.

A partir del análisis del mecanismo, se determinó que una de las principales causas de estas deficiencias correspondía a las características del sistema de guiado utilizado para el desplazamiento de los ejes. Esta condición generaba fricción, desalineaciones y movimientos irregulares durante la operación del brazo. Por ello, se planteó un rediseño mecánico parcial enfocado en mejorar el sistema de movimiento de los tres ejes mediante la incorporación de rodamientos lineales y nuevos ejes guía, con el propósito de aumentar la estabilidad estructural, reducir la resistencia al desplazamiento y mejorar la repetibilidad del sistema.

Además del rediseño mecánico, el trabajo incluyó la automatización del brazo robótico cartesiano para ampliar su funcionalidad dentro de una línea de producción. La aplicación seleccionada correspondió a una línea de empaquetado ubicada en el Laboratorio de Diseño de Procesos de la Universidad del Valle de Guatemala. Para ello, se integraron elementos de sensórica y control, entre ellos finales de carrera para la calibración y referencia de posición, una cámara para la implementación del sistema de visión computacional y dispositivos electrónicos para mejorar la interacción del sistema con su entorno. Asimismo, se emplearon microcontroladores para realizar las funciones de control, comunicación y coordinación del brazo robótico.

De esta manera, el proyecto dio continuidad al trabajo previo mediante la intervención de las limitaciones mecánicas identificadas y la incorporación de una arquitectura de automatización orientada a la ejecución de tareas de tipo \textit{pick and place}. Las modificaciones realizadas estuvieron dirigidas a obtener un funcionamiento más fluido, preciso y adecuado para su utilización dentro de un entorno de empaquetado controlado.